\documentclass[border=12pt]{standalone}
\usepackage[no-math]{fontspec}
\usepackage[brazil]{babel}
\usepackage{amsmath,amssymb}
\IfFileExists{C:/Windows/Fonts/Lato_Regular.ttf}{%
  \setmainfont{Lato}[
    Path = C:/Windows/Fonts/,
    UprightFont = Lato_Regular.ttf,
    BoldFont = Lato_Bold.ttf,
    ItalicFont = Lato_Italic.ttf,
    BoldItalicFont = Lato_BoldItalic.ttf
  ]
  \setsansfont{Lato}[
    Path = C:/Windows/Fonts/,
    UprightFont = Lato_Regular.ttf,
    BoldFont = Lato_Bold.ttf,
    ItalicFont = Lato_Italic.ttf,
    BoldItalicFont = Lato_BoldItalic.ttf
  ]
}{%
  \RequirePackage[default]{lato}
}
\renewcommand{\familydefault}{\sfdefault}

\usepackage{xcolor}
\usepackage{tikz}
\usepackage{pgfplots}
\pgfplotsset{compat=1.18}
\usetikzlibrary{arrows.meta,positioning,shapes.geometric,fit,backgrounds,calc,decorations.pathreplacing,patterns}

% Paleta institucional UEFS / Extensão Rural
\definecolor{erNavy}{HTML}{1F4268}
\definecolor{erNavyDark}{HTML}{152E4A}
\definecolor{erGold}{HTML}{B38A3D}
\definecolor{erGoldDark}{HTML}{8C6927}
\definecolor{erTeal}{HTML}{0E7C7B}
\definecolor{erForest}{HTML}{2D6A4F}
\definecolor{erAlert}{HTML}{9E2A2B}
\definecolor{erInk}{HTML}{26313D}
\definecolor{erMuted}{HTML}{5B6573}

\definecolor{bgNavy}{HTML}{EDF2F7}
\definecolor{bgGold}{HTML}{F8F3E9}
\definecolor{bgTeal}{HTML}{EBF6F5}
\definecolor{bgForest}{HTML}{EEF5F0}
\definecolor{bgAlert}{HTML}{FDF2F2}

\tikzset{
  >=Stealth,
  erBox/.style={
    draw=#1, fill=white, rounded corners=5pt, line width=1.15pt,
    inner sep=7pt, align=center, font=\footnotesize\sffamily, text=erInk
  },
  erTint/.style={
    draw=#1, fill=#1!8!white, rounded corners=5pt, line width=1.15pt,
    inner sep=7pt, align=center, font=\footnotesize\sffamily, text=erInk
  },
  erBadge/.style={
    fill=#1, text=white, rounded corners=9pt, inner xsep=8pt, inner ysep=3pt,
    font=\scriptsize\bfseries\sffamily
  },
  erTitle/.style={font=\normalsize\bfseries\sffamily, text=erNavy, align=center},
  erSub/.style={font=\footnotesize\sffamily, text=erMuted, align=center},
  erArrow/.style={->, >=Stealth, line width=1.35pt},
  erDArrow/.style={<->, >=Stealth, line width=1.35pt}
}


\begin{document}
\begin{tikzpicture}[node distance=1cm]

  % Título principal
  \node[erTitle, font=\large\bfseries\sffamily] (mainTitle) at (0, 5.2)
    {Evolução da Extensão Rural: Do Modelo Difusionista à PNATER Participativa};

  % ================= PAINEL 1: DIFUSIONISTA =================
  \begin{scope}[shift={(0, 2.55)}]
    % Fundo do painel
    \node[draw=erAlert!70, fill=bgAlert, rounded corners=8pt, line width=1.2pt,
          minimum width=17.2cm, minimum height=3.25cm] (panel1) at (0, 0) {};

    % Badge do painel
    \node[erBadge=erAlert, rounded corners=5pt, inner ysep=4pt, anchor=west]
      at ([xshift=14pt, yshift=-15pt]panel1.north west)
      {1. Modelo Difusionista Clássico (1950--1980)};
    \node[font=\scriptsize\itshape\sffamily, text=erAlert, anchor=east]
      at ([xshift=-14pt, yshift=-15pt]panel1.north east)
      {Fluxo unidirecional · Pacote tecnológico};

    % Nós do fluxo
    \node[erBox=erAlert, minimum width=3.3cm, minimum height=1.25cm] (p1_pesq) at (-6.0, -0.35)
      {\textbf{Pesquisador / ICT}\\[2pt]\scriptsize\color{erMuted}Gera tecnologia em estação};

    \node[erBox=erAlert, minimum width=3.3cm, minimum height=1.25cm] (p1_ext) at (0, -0.35)
      {\textbf{Extensionista}\\[2pt]\scriptsize\color{erMuted}Difusor / repassador técnico};

    \node[erBox=erAlert, minimum width=3.5cm, minimum height=1.25cm] (p1_agr) at (6.0, -0.35)
      {\textbf{Agricultor}\\[2pt]\scriptsize\color{erAlert}\textbf{(Receptor Passivo)}};

    % Setas unidirecionais com folga ampla
    \draw[erArrow, color=erAlert] (p1_pesq.east) -- (p1_ext.west)
      node[midway, above=5pt, font=\scriptsize\bfseries\sffamily, text=erAlert] {transfere pacote};
    \draw[erArrow, color=erAlert] (p1_ext.east) -- (p1_agr.west)
      node[midway, above=4pt, font=\scriptsize\bfseries\sffamily, text=erAlert, align=center]
      {transmite\\[-1pt]unidirecionalmente};
  \end{scope}

  % ================= PAINEL 2: DIALÓGICO & PNATER =================
  \begin{scope}[shift={(0, -1.85)}]
    % Fundo do painel
    \node[draw=erForest!80, fill=bgForest, rounded corners=8pt, line width=1.2pt,
          minimum width=17.2cm, minimum height=4.35cm] (panel2) at (0, 0) {};

    % Badge do painel
    \node[erBadge=erForest, rounded corners=5pt, inner ysep=4pt, anchor=west]
      at ([xshift=14pt, yshift=-15pt]panel2.north west)
      {2. Modelo Dialógico \& PNATER (Lei 12.188/2010)};
    \node[font=\scriptsize\itshape\sffamily, text=erForest, anchor=east]
      at ([xshift=-14pt, yshift=-15pt]panel2.north east)
      {Construção participativa · Agroecologia};

    % Nós do fluxo
    \node[erBox=erForest, minimum width=3.3cm, minimum height=1.25cm] (p2_pesq) at (-6.0, 0.2)
      {\textbf{Pesquisador}\\[2pt]\scriptsize\color{erMuted}Conhecimento científico};

    \node[erBox=erForest, minimum width=3.3cm, minimum height=1.25cm] (p2_ext) at (0, 0.2)
      {\textbf{Extensionista}\\[2pt]\scriptsize\color{erMuted}Educador e mediador};

    \node[erBox=erForest, minimum width=3.5cm, minimum height=1.25cm] (p2_agr) at (6.0, 0.2)
      {\textbf{Agricultor(a) Familiar}\\[2pt]\scriptsize\color{erForest}\textbf{(Sujeito e Saberes Locais)}};

    % Setas bidirecionais horizontais
    \draw[erDArrow, color=erForest] (p2_pesq.east) -- (p2_ext.west)
      node[midway, above=5pt, font=\scriptsize\bfseries\sffamily, text=erForest] {diálogo horizontal};
    \draw[erDArrow, color=erForest] (p2_ext.east) -- (p2_agr.west)
      node[midway, above=4pt, font=\scriptsize\bfseries\sffamily, text=erForest, align=center]
      {co-construção\\[-1pt]agroecológica};

    % Arco curvo de retroalimentação (Agricultor -> Pesquisador)
    \draw[erArrow, color=erTeal, line width=1.4pt, dashed]
      (p2_agr.south) .. controls +(0,-1.2) and +(0,-1.2) .. (p2_pesq.south)
      node[midway, fill=bgForest, inner xsep=10pt, inner ysep=2pt,
           font=\footnotesize\bfseries\itshape\sffamily, text=erTeal]
      {Diálogo de Saberes · Freire (retroalimentação e validação social)};
  \end{scope}

\end{tikzpicture}
\end{document}
